\chapter{Introduction}

\section{Background}
\label{sec:intro}

In sequential programs, state updates occur in a well-defined order, and each operation
observes a consistent state produced by previous operations. Concurrent programs break
this assumption by allowing multiple sequences of operations to interleave, which makes
it considerably harder to reason about what any individual operation will see.

Consider a simple example in which two threads concurrently increment a shared counter:

\begin{table}[ht]
  \centering
  \begin{tabular}{cc}
    \toprule
    \textbf{Thread 1} & \textbf{Thread 2} \\
    \midrule
    \texttt{ctr++}    & \texttt{ctr++}    \\
    \bottomrule
  \end{tabular}
\end{table}

\noindent At a high level, each increment appears atomic. On hardware, however,
it typically decomposes into three steps: reading the value from memory,
incrementing it, and writing it back:

\begin{table}[ht]
  \centering
  \begin{tabular}{cc}
    \toprule
    \textbf{Thread 1}        & \textbf{Thread 2}        \\
    \midrule
    $r_1 \gets \mathit{ctr}$ & $r_2 \gets \mathit{ctr}$ \\
    $r_1 \gets r_1 + 1$      & $r_2 \gets r_2 + 1$      \\
    $\mathit{ctr} \gets r_1$ & $\mathit{ctr} \gets r_2$ \\
    \bottomrule
  \end{tabular}
\end{table}

\noindent Since each thread performs three operations, the total number of possible
interleavings is $\displaystyle \binom{6}{3} = 20$, where only two are \emph{serializable}.

Correctness in concurrent systems therefore depends on controlling how shared
state is accessed and modified. Classical synchronization mechanisms achieve
this through mutual exclusion, ensuring that only one thread may access a
critical section at a time. Locks and monitors are among the most widely used
approaches for enforcing such guarantees.

While these mechanisms can provide thread safety, getting them right is
difficult in practice: locks must be manually acquired and released, and
preserving invariants in the presence of exceptions grows increasingly subtle
as systems scale. An exception thrown after a lock has been acquired but before
it is released can leave the lock held indefinitely, potentially deadlocking
the entire system.

Lomet~\cite{lomet1977} proposed a different framing: rather than managing
synchronisation explicitly, programmers should be able to declare that a block
of code must execute atomically and leave the mechanism to the underlying
system. Software Transactional Memory~(STM) is a realisation of this idea.

Beyond atomicity, STM provides something that lock-based approaches struggle
with: \emph{composability}. Consider a bank transfer — a withdrawal from one
account followed by a deposit into another. Each operation can be made atomic
in isolation, but composing them does not automatically make the pair atomic.
Another thread can observe the state between the two operations, seeing money
that has left one account but not yet arrived in the other.

\begin{listing}[ht]
  \begin{minted}{clojure}
(defn transfer [from to amount]
  (withdraw from amount)   ; atomic
  (deposit to amount))     ; atomic, but not together
\end{minted}
  \caption{Non-composable transfer}
\end{listing}

With STM, both operations are wrapped in a single transaction and committed
together:

\begin{listing}[ht]
  \begin{minted}{clojure}
(defn transfer [from to amount]
  (dosync
    (withdraw from amount)
    (deposit to amount)))
\end{minted}
  \caption{Composable transfer with STM}
\end{listing}

\noindent If a conflict occurs mid-execution, the runtime retries the entire
transaction from scratch, and no intermediate state is ever visible to other
threads. Each of \texttt{withdraw} and \texttt{deposit} can be written with
its own \texttt{dosync}; when called inside an enclosing transaction they
simply join it, which is the flat-nesting behaviour described in
Section~\ref{ch:implementation}.

\section{Problem Statement}

Clojure ships with its own STM implementation, built around \texttt{dosync}
and transactional references (\texttt{ref}s). It is well-designed and widely
used, but it has a documented limitation. The official documentation is direct
about it:

\begin{quote}
  I/O and other activities with side-effects should be avoided in
  transactions, since transactions will be retried~\cite{clojureRefs}.
\end{quote}

The reason is straightforward: if a transaction body is re-executed on conflict,
any side-effect inside it — a log write, a notification, an external API call —
will be repeated on every retry, not just on the successful commit. There is no
built-in mechanism to defer such effects to the point where the transaction is
known to have succeeded.

For some applications this is not a problem, but for others it is a real
constraint. Audit logging, for instance, needs to record exactly the commits
that took effect, not every attempt. Event-driven systems that fire
notifications on state changes face the same issue.

The most direct fix would be to extend Clojure's existing STM. The obstacle
is that the implementation lives in \texttt{LockingTransaction.java} inside
the \texttt{clojure.lang} package, and the methods that would need to be
extended — \texttt{doGet}, \texttt{doSet}, and \texttt{run} — have
package-private visibility. They cannot be overridden or wrapped from outside
the package. Jensen and Thomsen~\cite{jensen2016} worked around this by
forking the Clojure runtime itself and modifying \texttt{LockingTransaction}
directly, which works but ties the solution to a specific Clojure release and
adds ongoing maintenance burden.

This thesis takes a different route: rather than modifying the runtime, we
implement an STM system from scratch at the language level in Clojure. The
core algorithm is a variation of TL2~\cite{dice2006transactional}, extended
with MVCC and a transactional event system that exposes commit and rollback
hooks as first-class API. The result is decoupled from any particular Clojure
version and requires no changes to the runtime.

Table~\ref{tab:public-api} summarises the full public API. Transactional
operations live in the \texttt{opusdb.atomic.stm} namespace; the hook
mechanism is in \texttt{opusdb.atomic.tx-hooks}.

\begin{table}[H]
  \centering
  \renewcommand{\arraystretch}{1.4}
  \begin{tabularx}{\textwidth}{@{} l l X @{}}
    \toprule
    \textbf{Function / Macro} & \textbf{Namespace}                                                     & \textbf{Description} \\
    \midrule
    \texttt{(ref \textit{initial-val})}
                              & \texttt{stm}
                              & Creates a new transactional reference with the given initial value.                           \\
    \texttt{(deref \textit{r})}
                              & \texttt{stm}
                              & Reads reference \textit{r} within the current transaction, or directly
    if called outside of one.                                                                                                 \\
    \texttt{(ref-set \textit{r} \textit{v})}
                              & \texttt{stm}
                              & Buffers a write of value \textit{v} to reference \textit{r};
    must be called inside \texttt{dosync}.                                                                                    \\
    \texttt{(dosync \&\ \textit{body})}
                              & \texttt{stm}
                              & Executes \textit{body} atomically. Retries on conflict; composes with
    enclosing transactions via flat nesting.                                                                                  \\
    \texttt{(on-commit \textit{handler})}
                              & \texttt{tx-hooks}
                              & Registers \textit{handler} to be called exactly once after the
    enclosing transaction commits.                                                                                            \\
    \texttt{(on-rollback \textit{handler})}
                              & \texttt{tx-hooks}
                              & Registers \textit{handler} to be called on each abort of the
    enclosing transaction, before retry.                                                                                      \\
    \bottomrule
  \end{tabularx}
  \caption{Public API of the proposed STM system.}
  \label{tab:public-api}
\end{table}

\noindent A typical usage combining all six functions is shown in
Listing~\ref{lst:api-example}.

\begin{listing}[ht]
  \begin{minted}{clojure}
(dosync
  (let [bal (deref account)]
    (ref-set account (- bal amount))
    (on-commit  #(log/info "Transfer committed"))
    (on-rollback #(log/warn "Transfer retried"))))
\end{minted}
  \caption{Typical usage of the public API: transactional transfer with lifecycle hooks.}
  \label{lst:api-example}
\end{listing}